% =====================================================================
% SZKIC v1 (11.11.2026) -- wersja polska articles/jobshop/main.tex
% Do czytania i akceptacji; wersja do zgłoszenia to main.tex.
% WSZYSTKIE LICZBY W TABELACH 2-4 TO PILOT (1 s, prototyp w Pythonie,
% 2-3 ziarna): pokazują kierunek, zastąpi je kampania.
% Bibliografia: jak w main.tex (do sprawdzenia w Crossref).
% =====================================================================
\documentclass[10pt,twocolumn]{article}
\usepackage[margin=0.75in]{geometry}
\usepackage[T1]{fontenc}
\usepackage[utf8]{inputenc}
\usepackage[polish]{babel}
\usepackage{amsmath,amssymb,amsthm}
\usepackage{booktabs}
\usepackage{multirow}
\usepackage{xcolor}
\usepackage[hidelinks]{hyperref}

\newtheorem{theorem}{Twierdzenie}
\newtheorem{proposition}{Stwierdzenie}
\newtheorem{lemma}{Lemat}
\newtheorem{corollary}{Wniosek}
\newcommand{\todo}[1]{\textcolor{red}{[TODO: #1]}}
\newcommand{\pilot}{\textcolor{blue}{[pilot]}}
\newcommand{\JP}{\mathrm{JP}}
\newcommand{\JS}{\mathrm{JS}}
\newcommand{\Cmax}{C_{\max}}

\title{Złożone ruchy graniczne w problemie gniazdowym:\\
wykonalność przestawień w blokach\\
i kontinuum od flow shopu do job shopu}
\author{\todo{autorzy}}
\date{}

\begin{document}
\maketitle

\begin{abstract}
Sąsiedztwa złożone wykonują w jednym ruchu wiele niezależnych zamian,
a najlepszą kombinację wybiera programowanie dynamiczne. Sprawdzają się
w permutacyjnym problemie przepływowym, w którym każda permutacja jest
wykonalnym harmonogramem. Przenosimy je na problem gniazdowy z kryterium
$\Cmax$. Ścieżka krytyczna pełni rolę permutacji, a jej bloki rolę
odcinków, w których umieszcza się zamiany. Jedyną nową przeszkodą jest
wykonalność. Dowodzimy dokładnego warunku: przestawienie operacji
wewnątrz bloków krytycznych tworzy cykl wtedy i tylko wtedy, gdy odwraca
parę operacji połączoną ścieżką, która wychodzi z pierwszej przez jej
następnik technologiczny, wchodzi do drugiej przez jej poprzednik
technologiczny i omija resztę bloku. Z warunku wynika, że ruchy złożone
z rozłącznych przedziałów wymagają tylko testów pojedynczych składowych,
że krzyżujące się zamiany wprowadzają ograniczenia dla par, a w problemie
przepływowym żadne przestawienie nie jest niewykonalne. Klasyczna
własność blokowa sprowadza następnie każdą rodzinę zamian do najwyżej
dwóch ruchów na blok. Prowadzi to do rodziny ruchów granicznych o trzech
osiach: typ ruchu (zamiana albo wstawienie), zasięg oraz zastosowanie
pojedyncze albo złożone. N5, N6 i N7 są punktami tej rodziny, podobnie
jak złożone sąsiedztwa Fibonacciego, $k$-bonacciego i Dynasearch.
W badaniu pilotażowym z iterowanym przeszukiwaniem lokalnym ruchy złożone
o umiarkowanym zasięgu dają na instancjach Taillarda mniejszy $\Cmax$ niż
N5 i N6, a ich przewaga nad N5 jest największa dla instancji pośrednich
między flow shopem a job shopem.
\todo{ostatnie zdanie zastąpić wynikiem kampanii}
\end{abstract}

% ---------------------------------------------------------------------
\section{Wstęp}

Przeszukiwanie lokalne dla problemu gniazdowego z kryterium długości
uszeregowania, $J||\Cmax$, opiera się na kilku sąsiedztwach określonych na
ścieżce krytycznej. N1 zamienia dowolne dwie sąsiednie operacje krytyczne
jednej maszyny~\cite{laarhoven1992}. N5 zostawia tylko zamiany na obu
końcach każdego bloku krytycznego~\cite{nowicki1996}. N6 i N7 przenoszą
operację na początek albo na koniec jej
bloku~\cite{balas1998,zhang2007}. Każde z nich zmienia w jednym ruchu
jeden blok, a najlepsze algorytmy tabu~\cite{nowicki2005,peng2015} są
zbudowane właśnie na nich.

Dla permutacyjnego problemu przepływowego istnieje inna linia sąsiedztw.
Ruch złożony to zbiór zamian na rozłącznych przedziałach permutacji,
a najlepszy zbiór znajduje programowanie dynamiczne, jak
w Dynasearch~\cite{congram2002}. Ograniczenie odległości zamiany daje
sąsiedztwo Fibonacciego (tylko zamiany sąsiednie)~\cite{fibonacci_own}
i rodzinę $k$-bonacciego między Fibonaccim
a Dynasearch~\cite{kbonacci_own}. Sąsiedztwa te mają wykładniczy rozmiar
i przeszukuje się je tanio, ale korzystają z tego, że każda permutacja
jest wykonalnym harmonogramem przepływowym. W problemie gniazdowym
zamiana dwóch niesąsiednich operacji może zamknąć cykl w grafie
dysjunktywnym, a wykonalność kombinacji wykonalnych zamian nie jest
oczywista.

W tej pracy pytamy, co zostaje z sąsiedztw złożonych, gdy marszruty zadań
się różnią. Odpowiedź ma trzy części. Pierwsza to konstrukcja. Ścieżkę
krytyczną traktujemy jak wirtualną permutację, zamiana kandydująca leży
w jednym bloku, a wybór ruchu złożonego wykonuje ten sam program
dynamiczny co w problemie przepływowym. Program rozkłada się na bloki,
więc sąsiedztwo jest iloczynem sąsiedztw bloków. Druga część to
twierdzenie o wykonalności dowolnych przestawień wewnątrz bloków. Podaje
ono dokładny warunek dla par i wyjaśnia, które rodziny złożone są
bezpieczne: dla ruchów na przedziałach rozłącznych lub zagnieżdżonych
wystarczy sprawdzić każdą zamianę osobno, krzyżujące się zamiany wymagają
testów par, a przy wspólnych marszrutach nie trzeba sprawdzać niczego.
Trzecia część to uczciwy opis tego, co własność blokowa robi z tymi
rodzinami. Ruch, który nie zmienia żadnej granicy bloku, nie może
poprawić $\Cmax$, więc po zwykłym filtrze w każdym bloku zostają najwyżej
dwie zamiany, po jednej na każdym końcu. Rodziny liczone ciągami zapadają
się, a naturalną rodzinę dla problemu gniazdowego indeksują typ ruchu
granicznego, jego zasięg i to, czy stosujemy jeden ruch, czy po jednym na
każdy koniec bloku. N5, N6 i N7 są punktami tej rodziny z jednym ruchem.
Ich wersje złożone są nowe.

Część eksperymentalna mierzy trzy rzeczy w jednym protokole i jednym
kodzie: jakość osiąganą przez każdego członka rodziny w iterowanym
przeszukiwaniu lokalnym i symulowanym wyżarzaniu na instancjach
gniazdowych Taillarda~\cite{taillard1993}, udział niewykonalnych dalekich
ruchów w zależności od kształtu instancji i jakości rozwiązania oraz
zachowanie rodziny na generowanych instancjach pośrednich między flow
shopem a job shopem. \todo{jedno zdanie z głównym wynikiem kampanii}

% ---------------------------------------------------------------------
\section{Podstawy}

Instancja ma $n$ zadań i $m$ maszyn. Zadanie $j$ odwiedza maszyny
w kolejności swojej marszruty, a jego $t$-ta operacja trwa $p_{j,t}$. Dla
operacji $o$ przez $\JP(o)$ i $\JS(o)$ oznaczamy jej poprzednik
i następnik technologiczny. Rozwiązanie to kolejność $\sigma_k$ operacji
na każdej maszynie $k$. Graf dysjunktywny~\cite{roy1964} $G(\sigma)$ ma
łuk technologiczny $o\to\JS(o)$ dla każdej operacji i łuk maszynowy
między kolejnymi operacjami każdej $\sigma_k$. Rozwiązanie jest wykonalne
wtedy i tylko wtedy, gdy $G(\sigma)$ jest acykliczny. Głowa $r(o)$ to
długość najdłuższej ścieżki kończącej się na początku $o$, ogon $q(o)$ to
długość najdłuższej ścieżki wychodzącej z końca $o$,
a $\Cmax=\max_o\, r(o)+p_o+q(o)$. Jedno przejście w porządku
topologicznym liczy je wszystkie w czasie $O(nm)$.

Ścieżka krytyczna to najdłuższa ścieżka w $G(\sigma)$. Blok to maksymalna
seria kolejnych operacji krytycznych połączonych łukami maszynowymi.
Własność blokowa~\cite{brucker1994,nowicki1996} mówi, że w harmonogramie
o mniejszym $\Cmax$ pewna operacja bloku musi być wykonana przed pierwszą
albo po ostatniej operacji tego bloku. N5 korzysta z niej, zostawiając
dwie zamiany na blok. N6 i N7 definiują na jej podstawie wstawienia
i podają dla nich warunki wystarczające
wykonalności~\cite{balas1998,zhang2007}. Wykonalność wstawień w ogólnym
ujęciu badali Gr\"oflin i Klinkert~\cite{groflin2007}, a dla kryteriów
regularnych Mati i in.~\cite{mati2011}. \todo{po lekturze opisać
dokładnie relację twierdzenia~\ref{th:main} do
\cite{groflin2007,mati2011,dauzere1997} i do ruchów złożonych Bożejki
i Wodeckiego}
Kandydatów porządkuje się zwykle estymatą liczoną z bieżących głów
i ogonów~\cite{taillard1994,balas1998}, a dokładnie ocenia się tylko ruch
wybrany.

W permutacyjnym problemie przepływowym sąsiedztwa złożone liczą ciągi
liczbowe: permutacja długości $L$ ma $F_{L+1}$ zbiorów rozłącznych zamian
sąsiednich, $a_k(L)$ zbiorów rozłącznych zamian o odległości mniejszej niż
$k$ (liczby $k$-bonacciego) i $2^{L-1}$ zbiorów rozłącznych zamian
o dowolnej odległości. Dopuszczenie zamian zagnieżdżonych daje liczby
Motzkina, a dopuszczenie krzyżowań inwolucje~\cite{motzkin_own}.

% ---------------------------------------------------------------------
\section{Sąsiedztwa na ścieżce krytycznej}\label{sec:family}

Niech $P=(o_1,\dots,o_L)$ będzie jedną ścieżką krytyczną rozwiązania
wykonalnego, a $B_1,\dots,B_b$ jej blokami. Gdy oba poprzedniki operacji
są krytyczne, idziemy przez poprzednik maszynowy, co daje dłuższe bloki.
Pozycje $1,\dots,L$ ścieżki $P$ tworzą wirtualną permutację. Zamiana
wymienia dwie operacje jednego bloku i odpowiada przedziałowi pozycji.
Przedziały z różnych bloków są rozłączne, więc dowolny zbiór zamian
z różnych bloków jest ruchem złożonym w sensie problemu przepływowego.
Ruch przepływowy zmienia kolejność na wszystkich maszynach naraz. Ruch
gniazdowy zmienia jedną maszynę na blok.

\begin{proposition}[rozmiar]\label{pr:count}
Niech $a_F(L)$ będzie liczbą ruchów złożonych rodziny $F$ na permutacji
długości $L$. Na ścieżce krytycznej o długościach bloków $L_1,\dots,L_b$
rodzina ma $\prod_{i} a_F(L_i)$ ruchów, łącznie z pustym.
\end{proposition}

Iloczyn jest wykładniczy względem długości ścieżki, jak w problemie
przepływowym. Własność blokowa usuwa większość tych ruchów.

\begin{proposition}[reguła granic]\label{pr:border}
Niech $\sigma'$ będzie rozwiązaniem wykonalnym otrzymanym z $\sigma$
przez przestawienie operacji wewnątrz bloków ścieżki $P$. Jeśli $\sigma'$
zachowuje pierwszą operację każdego bloku poza $B_1$ i ostatnią operację
każdego bloku poza $B_b$, to $\Cmax(\sigma')\ge\Cmax(\sigma)$.
\end{proposition}
\begin{proof}
W $\sigma'$ operacje każdego bloku nadal stoją kolejno na swojej
maszynie. Odwiedzamy je w nowej kolejności. Łuk technologiczny z ostatniej
operacji $B_i$ do pierwszej operacji $B_{i+1}$ się nie zmienia, więc
odwiedzane operacje tworzą w $G(\sigma')$ ścieżkę o tych samych
operacjach co $P$ i tej samej długości.
\end{proof}

Zamiana, która nie dotyka żadnej granicy, ma efekt nieujemny, więc
program dynamiczny sumujący efekty zamian nigdy jej nie wybierze.
Zostawiamy zatem tylko kandydatów, którzy wymieniają pierwszą albo
ostatnią operację bloku. Filtr ten jest dokładny dla modelu addytywnego
używanego przy wyborze. Nie jest dokładny dla prawdziwego $\Cmax$ ruchu
złożonego, bo zamiana wewnętrzna może zmienić wynik zamiany granicznej
wykonanej razem z nią. Wracamy do tego w rozdziale~\ref{sec:eval}.

\begin{proposition}[spłaszczenie]\label{pr:flat}
Pod regułą granic blok długości $L$ wnosi do ruchu złożonego najwyżej
dwie rozłączne zamiany: jedną zawierającą jego pierwszą pozycję i jedną
zawierającą ostatnią. Zamiany zagnieżdżone nie występują, a dwie zamiany
krzyżują się tylko jako para $(1,j)$, $(i,L)$ dla $i<j$.
\end{proposition}

Stwierdzenie~\ref{pr:flat} zmienia obraz przeniesiony z problemu
przepływowego. Wewnątrz bloku rodzina Motzkina jest równa Dynasearch,
a liczba przefiltrowanych ruchów złożonych bloku jest wielomianowa
względem $L$. Zmieniać można już tylko sam ruch graniczny. Rozważamy trzy
typy ruchów na przedziale bloku $[a,b]$ zawierającym koniec bloku:
zamianę końców, wstawienie przenoszące operację wewnętrzną na koniec
bloku (N6) i wstawienie przesuwające operację końcową w głąb (N7).
Parametr zasięgu $k$ ogranicza przedział warunkiem $b-a\le k-1$. Dla
$k=2$ wszystkie trzy typy są tą samą zamianą sąsiednią. Sąsiedztwo
pojedyncze wykonuje najlepszego kandydata. Sąsiedztwo złożone wykonuje
najwyżej jednego kandydata na każdy koniec bloku we wszystkich blokach
naraz. Tabela~\ref{tab:family} wymienia badanych członków rodziny.

\begin{table}[t]
\centering\small
\begin{tabular}{@{}lll@{}}
\toprule
Typ i zasięg & Pojedyncze & Złożone \\
\midrule
zamiana, $k=2$ & N5 & Fib \\
zamiana, zasięg $k$ & & $k$-bonacci \\
zamiana, dowolny zasięg & & Dyn \\
wstawienie N6, dowolny zasięg & N6 & c6 \\
wstawienia N6 i N7, dowolny zasięg & N7 & c7 \\
wszystkie typy, zasięg $k$ & & bm$_k$ \\
\bottomrule
\end{tabular}
\caption{Rodzina ruchów granicznych. Pozycje c6, c7 i bm$_k$ są nowe.}
\label{tab:family}
\end{table}

Jedna iteracja ocenia bieżące rozwiązanie, wyznacza ścieżkę i jej bloki,
wylicza kandydatów granicznych wybranego typu i zasięgu, odrzuca
niewykonalnych (rozdział~\ref{sec:feas}), szacuje $\Cmax$ po każdym
z pozostałych, wybiera zbiór kandydatów i dokładnie ocenia otrzymane
rozwiązanie. Wybór wykonuje pozycyjny program dynamiczny z kodu
przepływowego. Kandydaci z różnych bloków nigdy się nie nakładają, więc
program zwraca sumę niezależnych optimów bloków i nie przedstawiamy go
jako wkładu. Jeśli nie wybierze niczego, wykonujemy najlepszego
pojedynczego kandydata, jak w naszych pracach o problemie przepływowym.
Koszt iteracji to dwie oceny $O(nm)$ i $O(bk^2)$ na estymaty dla rodziny
o zasięgu $k$.

% ---------------------------------------------------------------------
\section{Wykonalność przestawień w blokach}\label{sec:feas}

Niech $\sigma$ będzie wykonalne, a $\sigma'$ powstaje przez permutowanie
operacji wewnątrz bloków ścieżki $P$. Zmieniać się może kilka bloków
naraz. Dla dwóch operacji $x,y$ jednego bloku, gdzie $x$ stoi przed $y$
w $\sigma$, \emph{obejściem} z $x$ do $y$ nazywamy ścieżkę w $G(\sigma)$,
która wychodzi z $x$ łukiem technologicznym $x\to\JS(x)$, wchodzi do $y$
łukiem technologicznym $\JP(y)\to y$ i nie zawiera żadnej innej operacji
bloku. Dopuszczamy przypadek $\JS(x)=\JP(y)$.

\begin{lemma}\label{le:clean}
Obejście z $x$ do $y$ nie zawiera żadnej operacji ścieżki $P$ poza $x$
i $y$.
\end{lemma}
\begin{proof}
Niech $u$ będzie wewnętrzną operacją obejścia leżącą na $P$. Jeśli $u$
poprzedza $x$ na $P$, to $P$ prowadzi z $u$ do $x$, a obejście z $x$ do
$u$, co daje cykl w $G(\sigma)$. Jeśli $u$ następuje po $y$ na $P$, cykl
tworzą $P$ od $y$ do $u$ i obejście od $u$ do $y$. W pozostałym przypadku
$u$ leży na $P$ między $x$ a $y$, a więc należy do bloku, co wyklucza
definicja.
\end{proof}

\begin{theorem}\label{th:main}
Rozwiązanie $\sigma'$ jest niewykonalne wtedy i tylko wtedy, gdy istnieją
operacje $x,y$ jednego bloku takie, że $x$ poprzedza $y$ w $\sigma$, $y$
poprzedza $x$ w $\sigma'$, a $G(\sigma)$ zawiera obejście z $x$ do $y$.
\end{theorem}
\begin{proof}
Niech $\prec$ i $\prec'$ oznaczają porządki maszynowe w $\sigma$
i $\sigma'$. Różnią się one tylko na parach operacji tego samego bloku.

\emph{Dostateczność.} Z lematu~\ref{le:clean} wewnętrzne łuki obejścia
łączą operacje spoza $P$, a jego pierwszy i ostatni łuk są
technologiczne. Żaden z tych łuków nie zależy od kolejności wewnątrz
bloku, więc obejście jest ścieżką w $G(\sigma')$. Razem z $y\prec' x$
zamyka cykl.

\emph{Konieczność.} Niech $C$ będzie cyklem w $G(\sigma')$, przy czym
łuki maszynowe bierzemy w sensie przechodnim. Łuk maszynowy $u\to v$
cyklu $C$ nazywamy \emph{nowym}, jeśli $v\prec u$. Pozostałe łuki $C$
należą do $G(\sigma)$, więc $C$ zawiera nowy łuk i dzieli się na nowe
łuki oraz ścieżki z $G(\sigma)$ między nimi. Nowy łuk zostaje w jednym
bloku i cofa się po $P$. Ścieżka z $G(\sigma)$ między dwiema operacjami
$P$ nie może cofać się po $P$. Gdyby taka ścieżka przeszła z bloku do
późniejszego, $C$ nie mógłby wrócić, więc wszystkie nowe łuki cyklu leżą
w jednym bloku $B$.

Tniemy $C$ przy każdej wizycie w $B$. Fragment jest albo pojedynczym
łukiem maszynowym między operacjami $B$, zgodnym z $\prec'$, albo ścieżką
z $u\in B$ do $v\in B$, której wewnętrzne operacje leżą poza $B$.
Fragment drugiego rodzaju leży w $G(\sigma)$, więc $u\prec v$. Nie może
zaczynać się łukiem maszynowym, bo ten prowadziłby do operacji za blokiem
na tej samej maszynie, a reszta fragmentu wracałaby do $B$, zamykając
cykl w $G(\sigma)$. Z tego samego powodu nie może kończyć się łukiem
maszynowym. Jest więc obejściem z $u$ do $v$. Po scaleniu kolejnych
fragmentów maszynowych $C$ ma postać
$c_1\preceq' d_1\leadsto c_2\preceq' d_2\leadsto\dots\leadsto c_1$, gdzie
każde $\leadsto$ jest obejściem. Co najmniej jedno obejście występuje, bo
$\prec'$ jest porządkiem liniowym. Gdyby dla każdego obejścia zachodziło
$d_i\prec' c_{i+1}$, ciąg byłby cyklem w $\prec'$. Zatem pewne obejście
łączy parę, dla której $d_i\prec c_{i+1}$ i $c_{i+1}\prec' d_i$.
\end{proof}

Wykluczenie pozostałych operacji bloku jest istotne. Jeśli ścieżka z $x$
do $y$ może przechodzić przez operacje $z,z'$ bloku i używać łuku
maszynowego $z\to z'$, to nowa kolejność $z',y,x,z$ odwraca parę $(x,y)$,
a mimo to jest wykonalna. W testach pełnego przeglądu na losowych
instancjach warunek z twierdzenia~\ref{th:main} zgadzał się z wykrywaniem
cykli we wszystkich $184\,461$ sprawdzonych przestawieniach, z których
$28\,065$ było niewykonalnych, a wersja bez wykluczenia dała jeden
fałszywy alarm. \pilot

Twierdzenie zależy od $\sigma'$ tylko przez zbiór odwróconych par. Stąd
następujące wnioski.

\begin{corollary}\label{co:all}
(a) Zamiana dwóch sąsiednich operacji krytycznych jest
wykonalna~\cite{laarhoven1992}.
(b) Ruch złożony, którego składowe działają na przedziałach rozłącznych
lub zagnieżdżonych, jest wykonalny wtedy i tylko wtedy, gdy każda
składowa jest wykonalna osobno.
(c) Przeniesienie $u$ w przód za $x_1,\dots,x_s$ jest niewykonalne wtedy
i tylko wtedy, gdy istnieje obejście z $u$ do pewnego $x_i$.
Przeniesienie $u$ wstecz jest niewykonalne wtedy i tylko wtedy, gdy
istnieje obejście z pewnego $x_i$ do $u$.
(d) Jeśli wszystkie zadania mają tę samą marszrutę, każde przestawienie
wewnątrz bloków jest wykonalne.
\end{corollary}
\begin{proof}
(a) Obejście z $x$ do jego następnika maszynowego $y$ jest dłuższe niż
$p_x+p_y$, co przeczy krytyczności łuku $x\to y$. (b) Zbiór odwróconych
par takiego ruchu jest sumą zbiorów jego składowych. (c) Ruch odwraca
dokładnie pary $(u,x_i)$ albo $(x_i,u)$. (d) Przy wspólnej marszrucie
numer maszyny nigdy nie maleje wzdłuż ścieżki, a obejście najpierw go
zwiększa.
\end{proof}

Część (b) obejmuje złożone rodziny Fibonacciego, $k$-bonacciego,
Dynasearch i Motzkina oraz wszystkie złożone ruchy graniczne
z tabeli~\ref{tab:family}: ich wykonalność kosztuje jeden test na
kandydata i nic na kombinację. Krzyżujące się zamiany zachowują się
inaczej. Ich zbiór odwróconych par nie jest sumą składowych, dwie
wykonalne zamiany mogą dać ruch niewykonalny (0,71\% krzyżujących się par
w naszych testach \pilot), a zamianę niewykonalną może naprawić zamiana
krzyżująca się z nią. Każda odwrócona para zależy od najwyżej dwóch
zamian, więc wykonalność zbioru krzyżujących się zamian jest ograniczeniem
kwadratowym w zmiennych wyboru. Część (c) jest dokładna, podczas gdy
warunki używane z N6 i N7~\cite{balas1998,zhang2007} są wystarczające.
Część (d) wskazuje problem przepływowy jako przypadek bez przeszkód.

W implementacji parę sprawdza się najpierw głowami: ścieżka z $\JS(x)$
do $\JP(y)$ wymaga $r(\JP(y))\ge r(\JS(x))+p_{\JS(x)}$. Ten certyfikat
rozstrzyga około 90\% wykonalnych zamian niesąsiednich \pilot. Pozostałe
pary rozstrzyga przeszukiwanie w głąb z $\JS(x)$, przycinane głowami
i niewchodzące do bloku.

% ---------------------------------------------------------------------
\section{Ocena ruchu i addytywność}\label{sec:eval}

Dla kandydata, który przestawia przedział bloku, estymata przelicza
w nowej kolejności głowy operacji przedziału od lewej do prawej i ich
ogony od prawej do lewej. Głowy poprzedników technologicznych
i maszynowego poprzednika przedziału oraz ogony następników
technologicznych i maszynowego następnika pochodzą z bieżącego
rozwiązania. Estymatą nowego $\Cmax$ jest największa wartość $r'+p+q'$
w przedziale, a jej koszt jest liniowy względem długości przedziału. Dla
zamiany sąsiedniej to estymata Taillarda~\cite{taillard1994}: dolne
ograniczenie nowego $\Cmax$, dokładne zawsze wtedy, gdy ruch nie
poprawia. Dla dłuższych przedziałów to przybliżenie. \todo{E0: korelacja
rang z wartością dokładną i udział fałszywych popraw, dla każdego
zasięgu}

Wybór sumuje oszacowane efekty wybranych kandydatów. W problemie
gniazdowym suma ta zawyża rzeczywisty zysk bardziej niż w przepływowym,
bo skrócenie jednego bloku czyni krytyczną ścieżkę równoległą. W pilocie
rzeczywisty zysk ruchu złożonego wynosił 36--43\% sumy dokładnych efektów
pojedynczych i 14--19\% sumy estymat, a mimo to tylko 2--4\% ruchów
złożonych nie poprawiało. \pilot{} Dokładna ocena wybranego ruchu jest
więc konieczna. \todo{E5: środki zaradcze, czyli dokładna ocena $t$
najlepszych ruchów, limit liczby składowych oraz wariant bez filtra, który
mierzy, ile traci filtr granic}

% ---------------------------------------------------------------------
\section{Protokół eksperymentów}\label{sec:setup}

Protokół jest taki jak w naszych pracach o problemie przepływowym, żeby
jedynym zmienianym czynnikiem było sąsiedztwo. Iterowane przeszukiwanie
lokalne i symulowane wyżarzanie zachowują przebieg sterowania i stałe.
\todo{stałe w jednym zdaniu, jak w pracy o $k$-bonaccim} Rozwiązanie
początkowe przydziela operacje zadaniom po kolei i dekoduje je
półaktywnie, co odpowiada permutacji identycznościowej. Zaburzeniem jest
double bridge na porządku topologicznym operacji z tym samym
dekodowaniem, więc zawsze daje rozwiązanie wykonalne. Identyfikatory tabu
to zbiory zamienionych par operacji. Wszystkie sąsiedztwa mają jedną
implementację, która różni się tylko generatorem kandydatów.

Głównym zestawem testowym są instancje ta01--ta80
Taillarda~\cite{taillard1993}, osiem klas od $15\times15$ do
$100\times20$ po dziesięć instancji. Wszystkie instancje odtworzyliśmy
z opublikowanych ziaren. Każdą konfigurację uruchamiamy raz na 10~s
z pięcioma ziarnami, a wyniki dla budżetów 0,1, 0,5, 1, 2 i 5~s
odczytujemy z dziennika ruchów tego przebiegu. Podajemy względne
odchylenie od najlepszego znanego górnego
ograniczenia~\cite{vanhoorn2018} z datą i źródłem każdej wartości oraz
liczbę iteracji przy każdym budżecie czasu. \todo{maszyna, wersja
Pythona; walidacja ekstrakcji budżetów na dwóch klasach}

Eksperyment kontinuum używa generatora Taillarda z parametrem nieporządku
marszrut $\rho\in\{0;0{,}1;0{,}25;0{,}5;1\}$. Marszruty zaczynają jako
identyczność, a każdą transpozycję tasowania marszruty wykonuje się
z prawdopodobieństwem $\rho$. Czasy operacji nie zależą od $\rho$.
Wartość $\rho=0$ daje problem przepływowy, w którym kolejności na
maszynach mogą się różnić, a $\rho=1$ standardową instancję gniazdową.
Rozmiary to $20\times5$, $20\times10$, $30\times15$ i $50\times10$ po
dziesięć instancji. Jako most do literatury przepływowej instancje
przepływowe Taillarda o rozmiarach $20\times5$, $20\times10$
i $50\times10$ rozwiązujemy także jako gniazdowe o wspólnych marszrutach.

Warianty porównujemy testem Friedmana z różnicą krytyczną
Nemenyiego~\cite{demsar2006} po instancjach, a kluczowe pary testem
Wilcoxona z poprawką Holma przy $\alpha=0{,}05$.

% ---------------------------------------------------------------------
\section{Wyniki}\label{sec:exp}

\todo{Ten rozdział zawiera tylko wartości pilotażowe. Docelowy układ:
zejście do optimów lokalnych; anatomia ścieżki i wykonalność; główne
porównanie na ta01--ta80 z diagramem różnic krytycznych; kontinuum;
addytywność; odniesienie do przeszukiwania tabu.}

\paragraph{Anatomia.}
Tabela~\ref{tab:anatomy} opisuje ścieżkę krytyczną rozwiązań osiągniętych
krótkim iterowanym przeszukiwaniem lokalnym. Bloki wydłużają się ze
wzrostem $n/m$. Przy małym zasięgu regułę granic przechodzi niewielu
kandydatów, a w klasach z $n/m\ge2$ duża część dalekich zamian jest
niewykonalna. Użyteczny zasięg ogranicza więc wykonalność, a nie tylko
koszt oceny.

\begin{table}[t]
\centering\small
\begin{tabular}{@{}lrrrrr@{}}
\toprule
Klasa & bloki & długość & $k{=}4$ & $k{=}\infty$ & niewyk. \\
\midrule
$15\times15$  & 7,7  & 3,1  & 12,2 & 18,3  & 19\% \\
$20\times20$  & 13,0 & 2,4  & 18,7 & 28,2  & 20\% \\
$30\times15$  & 4,7  & 9,1  & 14,0 & 41,8  & 45\% \\
$50\times15$  & 3,7  & 20,1 & 11,7 & 63,7  & 57\% \\
$100\times20$ & 3,7  & 32,8 & 12,3 & 116,7 & 67\% \\
\bottomrule
\end{tabular}
\caption{Bloki, średnia długość bloku, kandydaci zamian dla zasięgu 4
i nieograniczonego oraz udział niewykonalnych wśród tych ostatnich.
\pilot}
\label{tab:anatomy}
\end{table}

\paragraph{Rodzina ruchów granicznych.}
Tabela~\ref{tab:ils} porównuje dziewięciu członków rodziny w iterowanym
przeszukiwaniu lokalnym. Złożone zamiany o zasięgu 3 były lepsze od N5 na
wszystkich jedenastu instancjach i lepsze od N6 na dziewięciu. Złożona
wersja N7 była lepsza od N7 na dziesięciu instancjach. Dla N6 wersja
złożona wygrała na siedmiu instancjach i przegrała na trzech, czego test
znaków nie rozróżnia. Warianty złożone są blisko siebie. Pilot sugeruje
więc, że wykonanie ruchu granicznego we wszystkich blokach naraz znaczy
więcej niż typ ruchu, a dla N6 zostawia pytanie otwarte.

\begin{table}[t]
\centering\small
\begin{tabular}{@{}lrrr@{}}
\toprule
Wariant & RPD [\%] & śr. ranga & iter./s \\
\midrule
N5        & 10,62 & 8,77 & 2868 \\
N6        & 8,41  & 5,73 & 2304 \\
N7        & 8,90  & 7,05 & 1970 \\
$k=3$     & 7,40  & 3,36 & 2659 \\
Dyn       & 7,47  & 4,45 & 2095 \\
c6        & 7,53  & 3,68 & 2109 \\
c7        & 7,50  & 4,32 & 1974 \\
bm$_3$    & 7,42  & 3,68 & 2403 \\
bm$_\infty$ & 7,70 & 3,95 & 1592 \\
\bottomrule
\end{tabular}
\caption{Iterowane przeszukiwanie lokalne, 1~s, jedenaście instancji,
trzy ziarna; odchylenie od najlepszego znanego $\Cmax$. \pilot}
\label{tab:ils}
\end{table}

\paragraph{Od flow shopu do job shopu.}
Tabela~\ref{tab:rho} pokazuje dwa większe rozmiary eksperymentu
kontinuum. Przy wspólnych marszrutach N5 dało najlepszy wynik dla każdego
rozmiaru. Dla $\rho$ od 0,1 do 0,5 traciło do najlepszego wariantu od 14
do 22 punktów procentowych, a warianty złożone różniły się między sobą
o kilka punktów. Przy $\rho=1$ różnica malała, a dla $n/m\ge4$ znikała,
bo ścieżka krytyczna dobrego rozwiązania jest jednym blokiem, a $\Cmax$
równa się ograniczeniu maszynowemu. Przed interpretacją tego wzoru trzeba
rozdzielić dwa efekty. Przy $\rho=0$ N5 wykonuje od dwóch do czterech
razy więcej iteracji niż warianty złożone, a łagodny nieporządek marszrut
wydłuża bloki. \todo{etap continuum: wyniki przy równej liczbie iteracji,
długości bloków i udziały niewykonalnych dla każdego $\rho$}

\begin{table}[t]
\centering\small
\begin{tabular}{@{}lrrrrr@{}}
\toprule
Rozmiar & $\rho$ & N5 & $k{=}3$ & Dyn & c6 \\
\midrule
\multirow{5}{*}{$30\times15$}
 & 0    & 0,68  & 2,46 & 2,27 & 2,47 \\
 & 0,1  & 18,74 & 3,39 & 7,52 & 4,13 \\
 & 0,25 & 16,62 & 4,02 & 4,26 & 2,20 \\
 & 0,5  & 14,18 & 3,85 & 4,28 & 2,61 \\
 & 1    & 5,78  & 2,81 & 1,84 & 1,59 \\
\midrule
\multirow{5}{*}{$50\times10$}
 & 0    & 0,64  & 1,41 & 1,65 & 1,11 \\
 & 0,1  & 22,14 & 7,54 & 4,37 & 5,63 \\
 & 0,25 & 17,75 & 4,58 & 2,69 & 4,39 \\
 & 0,5  & 17,35 & 5,72 & 1,75 & 2,63 \\
 & 1    & 1,59  & 0,48 & 0,00 & 0,00 \\
\bottomrule
\end{tabular}
\caption{Odchylenie [\%] od najlepszego wyniku dowolnego wariantu na tej
samej instancji dla nieporządku marszrut $\rho$. \pilot}
\label{tab:rho}
\end{table}

% ---------------------------------------------------------------------
\section{Uwaga o sformułowaniach kwantowych}

Wybór ruchu złożonego jest kwadratowym problemem binarnym bez ograniczeń
z jedną zmienną na kandydata i karą za każdą parę nakładających się
kandydatów, jak w naszej pracy o problemie
przepływowym~\cite{windowed_own}. W problemie gniazdowym zadanie to jest
blokowo-diagonalne i małe: blok wnosi najwyżej $2(k-1)$ zmiennych,
a całe ruchy miały w dobrych rozwiązaniach od 4 do 19 zmiennych dla
$k\le4$. \pilot{} Cały ruch mieści się na obecnym urządzeniu bramkowym
bez okien. Program dynamiczny rozwiązuje to samo zadanie dokładnie, więc
nie twierdzimy, że daje to przyspieszenie. Ciekawym przypadkiem są rodziny
z krzyżowaniami, bo ich ograniczenia wykonalności są prawdziwymi wyrazami
dwuciałowymi. \todo{zdecydować, czy ten rozdział zostaje, czy idzie do
osobnej pracy}

% ---------------------------------------------------------------------
\section{Wnioski}

\todo{Napisać po kampanii. Przewidywana treść: zalecany zasięg
w zależności od $n/m$; rola złożenia wobec typu ruchu; w którym miejscu
osi od flow shopu do job shopu ruchy złożone się opłacają; ograniczenia,
czyli szkielet iterowanego przeszukiwania lokalnego, implementacja
w Pythonie i kryterium $\Cmax$.}

\todo{Bibliografia jak w main.tex.}
\begin{thebibliography}{99}\small
\bibitem{balas1998} Balas, Vazacopoulos, Manag. Sci. 1998.
\bibitem{brucker1994} Brucker, Jurisch, Sievers, Discrete Appl. Math. 1994.
\bibitem{congram2002} Congram, Potts, van de Velde, INFORMS J. Comput. 2002.
\bibitem{dauzere1997} Dauz\`ere-P\'er\`es, Paulli, Ann. Oper. Res. 1997.
\bibitem{demsar2006} Dem\v{s}ar, J. Mach. Learn. Res. 2006.
\bibitem{fibonacci_own} praca własna: sąsiedztwo Fibonacciego.
\bibitem{groflin2007} Gr\"oflin, Klinkert, Eur. J. Oper. Res. 2007.
\bibitem{kbonacci_own} praca własna: rodzina $k$-bonacciego.
\bibitem{laarhoven1992} van Laarhoven, Aarts, Lenstra, Oper. Res. 1992.
\bibitem{mati2011} Mati, Dauz\`ere-P\'er\`es, Lahlou, Eur. J. Oper. Res. 2011.
\bibitem{motzkin_own} praca własna: sąsiedztwo Motzkina.
\bibitem{nowicki1996} Nowicki, Smutnicki, Manag. Sci. 1996.
\bibitem{nowicki2005} Nowicki, Smutnicki, J. Sched. 2005.
\bibitem{peng2015} Peng, L\"u, Cheng, Comput. Oper. Res. 2015.
\bibitem{roy1964} Roy, Sussmann, SEMA 1964.
\bibitem{taillard1993} Taillard, Eur. J. Oper. Res. 1993.
\bibitem{taillard1994} Taillard, ORSA J. Comput. 1994.
\bibitem{vanhoorn2018} van Hoorn, J. Sched. 2018.
\bibitem{windowed_own} praca własna: okienkowe sąsiedztwa QUBO.
\bibitem{zhang2007} Zhang, Li, Guan, Rao, Comput. Oper. Res. 2007.
\end{thebibliography}

\end{document}
